\documentclass[11pt,a4paper]{article}
\usepackage[T1]{fontenc}
\usepackage{microtype}
\usepackage{isabelle,isabellesym}
\usepackage{amsmath}
\usepackage{amssymb}
\usepackage{url}
\usepackage{booktabs}
\usepackage{tabularx}
\usepackage{seqsplit}

\newcolumntype{Y}{>{\raggedright\arraybackslash}X}
% \thy / \thm / \cmd render Isabelle identifiers in typewriter and allow
% line breaks at every underscore.  The url package's \path command
% already supports this and handles the literal underscore character
% without needing math-mode escaping.  When invoked inside math mode
% we wrap in \mbox{} so \path's verbatim-like catcode changes do not
% corrupt the surrounding math context.
\newcommand{\bmsspidpath}[1]{\path{#1}}
\newcommand{\bmsspid}[1]{\ifmmode\mbox{\bmsspidpath{#1}}\else\bmsspidpath{#1}\fi}
\newcommand{\thy}[1]{\bmsspid{#1}}
\newcommand{\thm}[1]{\bmsspid{#1}}
\newcommand{\cmd}[1]{\bmsspid{#1}}
\emergencystretch=3em

\isabellestyle{it}
% Isabelle's italic document style renders underscores as hyphens in
% generated proof text.  Make those separators breakable so long theorem
% and rule names do not run past the page edge.
\def\isacharunderscore{\discretionary{-}{}{-}}
\def\isacharunderscorekeyword{\discretionary{-}{}{-}}

\begin{document}

\title{Correctness and Complexity of BMSSP in Isabelle/HOL}
\author{Arthur Freitas Ramos \and David Barros Hulak \and Ruy J. G. B. de Queiroz}
\maketitle

\begin{abstract}
This development formalizes the bounded multi-source shortest path (BMSSP)
algorithm underlying the deterministic directed single-source shortest path
algorithm of Duan, Mao, Mao, Shu, and Yin~\cite{DuanMaoMaoShuYin2025}.
It proves five explicit claims.  First, the executable
\texttt{bmssp\_distances} computes the exact reachable shortest-distance map
for well-formed finite natural weighted digraphs.  Second, a costed
BMSSP recurrence satisfies the \(O(m \log^{2/3} n)\) bound under the
abstract Insert/BatchPrepend/Pull operation-cost interface.  Third, the
bucketed partition operations realize the primitive costs required by that
interface.  Fourth, the costed recurrence and the correctness statement are
joined on a single object: the running time bounded by the headline is, for
all but finitely many sizes, the cost of a concrete BMSSP run that provably
returns the exact single-source shortest-path distances on the whole vertex
set.  Fifth, for the unbounded family of unit-weight directed paths \(P_k\)
(with \(\lvert E\rvert = k\) and \(\lvert V\rvert = k+1\) both growing with
\(k\)), the closed BMSSP cost \emph{expression} is proved
\emph{unconditionally} to be \(O(k \log^{2/3} k)\) in the true graph size
\(k\).  We are deliberately explicit about a limitation we uncovered: the
companion \emph{algorithm-level} size-parametric statement carries a
run-existence hypothesis that is unsatisfiable for \(P_k\) at \(k \ge 2\).  We
machine-check the structural core of this obstruction --- a sound exact-concrete
loop is forced into at most one genuine step --- from which the impossibility
follows, since that single step cannot cover the whole path; the statement,
though valid, is thus vacuous on its own family and is retained only as a
disclosed obstruction.  The
formalization does not state a single generated-code machine-step theorem for
the executable.  The entry also exports executable
SML and includes a checked worked shortest-path example whose evaluated
output is \([(0,0),(1,3),(2,5),(3,8),(4,6)]\).

AI assistance was used for proof engineering. The final definitions,
statements, and proofs are checked by Isabelle.
\end{abstract}

\tableofcontents

\section{Background}

The single-source shortest path problem asks for the distance from one source
vertex to every reachable vertex of a directed graph with non-negative edge
weights.  Dijkstra's algorithm~\cite{Dijkstra1959} is the classical starting
point: it maintains tentative labels and repeatedly settles a vertex of
minimum tentative distance.  With Fibonacci heaps, Fredman and
Tarjan~\cite{FredmanTarjan1987} obtained the well-known
\(O(m+n\log n)\) implementation for graphs with \(n\) vertices and \(m\)
edges.  For sparse directed graphs this bound became the standard deterministic
barrier: the logarithmic sorting-like term is paid by the priority-queue
discipline that extracts vertices one by one.  The barrier had previously been
broken only under additional assumptions; most notably,
Thorup~\cite{Thorup1999} solves the \emph{undirected} problem with positive
integer weights in linear time using a hierarchical bucketing structure.  The
result of Duan, Mao, Mao, Shu, and Yin is a comparison-based improvement for
the harder \emph{directed} problem.

Duan, Mao, Mao, Shu, and Yin~\cite{DuanMaoMaoShuYin2025} break this barrier
by changing the recursive unit of work.  Their algorithm is not simply
Dijkstra with a faster heap.  The central object is BMSSP, a bounded
multi-source shortest path subproblem.  A BMSSP call starts from a set of
already relevant sources and an upper bound \(B\).  It only has to complete
vertices whose shortest distances fall below the bound and whose shortest
paths pass through the current source set.  Recursive calls therefore work on
ranges of distances rather than on one global queue of all unsettled vertices.

The formalization follows that structure.  At the correctness level, the
algorithm is expressed in terms of finite directed graphs, simple walks,
shortest-path distances, and BMSSP pre/postconditions.  At the cost level, the
recursive proof is parameterized by an abstract partition interface with
Insert, BatchPrepend, and Pull operations.  The runtime theorem is proved once
against that interface and the paper's logarithmic schedule.  The concrete
bucketed data structure is then shown to realize the interface with the
critical operation bounds.  This cost statement is separate from the
functional-correctness theorem for the exported executable.

The essential improvement is the ratio inside the logarithm.  A sorted-list
or ordinary balanced-tree representation searches among all \(N\) live
entries and naturally gives a \(\log N\) term.  The paper's partition
structure searches among bucket boundaries, so Insert pays
\(O(\max(1,\log(N/M)))\), where \(M\) is the block-size parameter.  With the
schedule used by the recursive BMSSP analysis, this distinction is exactly
what preserves the \(O(m \log^{2/3} n)\) headline.

\section{Roadmap of the Formalization}

The formal development is organized by proof layer.  The table below gives one
or two primary entry points for each layer and the checked surface that a
reader should look for there; supporting theories appear nearby in the session
order.

\begin{center}
\small
\begin{tabularx}{\textwidth}{@{}>{\raggedright\arraybackslash}p{0.20\textwidth}
  >{\raggedright\arraybackslash}p{0.28\textwidth}Y@{}}
\toprule
\textbf{Layer} & \textbf{Start here} & \textbf{Checked surface} \\
\midrule
\textbf{Semantic target} &
\thy{BMSSP_Correctness} &
Finite weighted digraphs, shortest-path distances, BMSSP pre/postconditions,
and root-call correctness for ordinary SSSP. \\
\addlinespace
\textbf{Recursive BMSSP} &
\thy{BMSSP_Algorithm_Correctness};
\thy{BMSSP_Recursive} &
Correctness of FindPivots, base cases, partition-loop traces, recursive calls,
and final assembly of completed vertices. \\
\addlinespace
\textbf{Partition interface} &
\thy{BMSSP_Partition_Interface};
\thy{BMSSP_Partition_Pull_Bridge} &
Abstract Insert, BatchPrepend, and Pull contracts, with the algebraic cost
predicates used by the recurrence proof. \\
\addlinespace
\textbf{Cost recurrence} &
\thy{BMSSP_Top_Level_Bounds} &
The logarithmic parameter schedule and the cost-model theorem
\thm{bmssp_runtime_bigo_target}. \\
\addlinespace
\textbf{Verified-runtime bridge} &
\thy{BMSSP_Verified_Runtime} &
Couples the asymptotic bound with correctness on one object:
\thm{verified_bmssp_runtime_bigo_target} shows the bounded running time is,
eventually, the cost of a run proved \thm{sssp_correct} on the whole vertex
set. \\
\addlinespace
\textbf{Size-parametric family} &
\thy{BMSSP_Path_Family} &
For the growing path family \(P_k\), \thm{path_size_cost_bigo_size}
(unconditional) proves the closed cost \emph{expression} is \(O(k \log^{2/3} k)\)
in the true graph size \(k\); \thm{path_k_bounded_instance} certifies locale
inhabitation for every \(k\). The algorithm-level companion
\thm{bmssp_path_family_runtime_bigo_size} is conditional and, for this family,
vacuous (see the obstruction row).  The checked charged bridge
\thm{bmssp_path_family_charged_direct_insert_runtime_bigo_size_if_cost_bounded}
records the strong-version integration point, but still assumes eventual charged
run existence and the charged runtime chain ending in
\thm{charged_direct_insert_top_level_cost_refined_bound_log_params_fixed_degree},
which remains conditional on
\thm{charged_direct_insert_closed_refined_bound_log_params_fixed_degree}. \\
\addlinespace
\textbf{Totality obstruction} &
\thy{BMSSP_Cost_Totality} &
Machine-checked record that a sound exact-concrete loop is forced into at most
one genuine step, and none at a strictly positive lower bound
(\thm{exact_concrete_step_forces_empty_prefix},
\thm{no_genuine_exact_concrete_step_at_positive_lower_bound}); the explicit
reduction recorded there shows this single step cannot settle the whole path,
so no complete run exists on \(P_k\) for \(k \ge 2\) and the \texttt{runs\_exist}
hypothesis of the algorithm-level size-parametric bound is unsatisfiable for the
path family. \\
\addlinespace
\textbf{Bucketed data structure} &
\thy{BMSSP_Bucketed_Partition} &
Concrete bucketed operations satisfying
\thm{bp_insert_cost_bound}, \thm{bp_batch_prepend_cost_bound},
\thm{bp_pull_cost_bound}, and the interface facts
\thm{bp_realises_partition_insert_cost_bound},
\thm{bp_realises_partition_batch_cost_bound}, and
\thm{bp_realises_partition_pull_cost_bound}. \\
\addlinespace
\textbf{Executable artifact} &
\thy{BMSSP_Executable_Headline} &
The executable \texttt{bmssp\_distances}, the checked example
\thm{example_bmssp_correct}, and the headline theorem
\thm{bmssp_correct_strong}. \\
\bottomrule
\end{tabularx}
\end{center}

Several additional theories separate the mathematical layers used below:
shortest-path semantics, abstract BMSSP correctness, costed recurrences,
partition-interface costs, bucketed primitive operations, executable code, and
the final executable wrapper.  The cost-model theorems and the executable
correctness theorem are deliberately separate claims.

\section{Assumptions and Scope}

The entry proves three related, but not identical, layers.

First, the semantic BMSSP development starts from finite weighted digraphs,
simple walks, shortest-path distances, and a bounded BMSSP pre/postcondition.
The recursive correctness proof is organized through FindPivots, partition
loop obligations, complete sources, and bounded distance ranges.

Second, the asymptotic development is a cost-model proof.  It solves the
costed BMSSP recurrence under the abstract Insert/BatchPrepend/Pull interface,
uses the logarithmic parameter schedule from the paper, and proves the
\thm{bmssp_runtime_bigo_target} theorem in the corresponding runtime locale.
This is not stated as a machine-step bound for generated SML.

Third, the executable development proves functional correctness for the
exported finite natural-weight graph representation.  The public theorem
\thm{bmssp_correct_strong} applies to every \(G\) satisfying
\thm{nat_graph_well_formed}\,\(G\) and every source
\(\text{src} \in \thm{nat_graph_vertices}\,G\).  It does not assume unique
shortest paths.  The well-formedness predicate is a representation-side
condition: the directed edge list is duplicate-free and the total natural edge
weight is below the executable sentinel \thm{bmssp_infinity}.

\section{The Headline Theorems}
\label{sec:headline}

The primary executable correctness statement is

\begin{center}
\thm{bmssp_correct_strong} \quad in \quad \thy{BMSSP_Executable_Headline}.
\end{center}

The theorem fixes a finite directed graph \(G\) of natural-numbered vertices
with non-negative natural edge weights and a source vertex \(\text{src}\).
The statement is

\begin{quote}
For every \(G\) with \thm{nat_graph_well_formed}\,\(G\) and every
\(\text{src}\) in \thm{nat_graph_vertices}\,\(G\):
\begin{itemize}
\item the keys of \(\text{bmssp\_distances}\,G\,\text{src}\) are distinct;
\item the key set is exactly the set of vertices reachable from \(\text{src}\);
\item every pair \((v, d)\) returned by \(\text{bmssp\_distances}\,G\,\text{src}\)
      satisfies \(d = \text{nat\_graph\_dist}\,G\,\text{src}\,v\)
      after coercion to real numbers.
\end{itemize}
\end{quote}

Thus the executable theorem is not conditional on a unique-shortest-path or
tie-breaking assumption.  Its displayed hypotheses are the finite
natural-graph representation condition and the source-membership condition.

The runtime and bucketed-interface statements are separate public claims.
\thm{bmssp_runtime_bigo_target} is the cost-model theorem for the abstract
BMSSP recurrence and logarithmic parameter schedule.  Its closed
assumption-free form is stated inside the runtime locale; interpreting that
locale supplies the reduced positive instance, reachability, and bounded
outdegree hypotheses.  To certify that this locale is inhabited rather than
vacuous, the theory \thy{BMSSP_Runtime_Instance} interprets it on a concrete
unit-weight directed path \(0 \to 1 \to 2 \to 3\): it discharges unique
shortest walks (from the verified simple-walk enumerator), positive weights,
full reachability, and bounded outdegree, and thereby obtains
\thm{path_runtime_bigo_target}, the same logarithmic running-time bound as a
closed, assumption-free statement for an explicit non-trivial graph.  The
corollaries
\thm{bp_realises_partition_insert_cost_bound},
\thm{bp_realises_partition_batch_cost_bound}, and
\thm{bp_realises_partition_pull_cost_bound} show that the bucketed Insert,
BatchPrepend, and Pull operations realize the abstract primitive-cost
predicates used by that recurrence.  The more concrete theorems
\thm{bp_insert_cost_bound}, \thm{bp_batch_prepend_cost_bound}, and
\thm{bp_pull_cost_bound} expose the amortised budgets for the bucketed
implementation itself.

The session build checks all of these claims.  The executable smoke test
\thm{example_bmssp_correct} also executes the generated code equations on a
worked five-vertex example with
\(\text{bmssp\_distances}\,\text{example\_graph}\,0 =
  [(0,0),(1,3),(2,5),(3,8),(4,6)]\), proved by \texttt{by eval} inside the
session build itself.

\section{The Bucketed Partition Data Structure}

The partition data structure is the main concrete contribution of the
executable layer.  It stores tentative labels as key/value pairs grouped into
buckets.  Each bucket contains an unsorted list of at most \(M\) entries in
the regular state, and the bucket markers form a monotone directory of lower
bounds.  The directory, rather than the full set of entries, is the object
searched by Insert.  If \(N\) entries have been inserted since the last Pull,
there are only on the order of \(N/M\) buckets, which gives the
\(\log(N/M)\) search term.

The formal model separates several invariants.  The ordinary invariant states
that bucket sizes are within the block-size bound, markers are sorted, markers
are lower bounds for their bucket entries, keys are distinct, and the
functional value memory agrees with the entries stored in buckets.  The
ordered invariant additionally states that adjacent bucket boundaries separate
the value ranges.  The lazy invariant relaxes the size side condition: a
bucket may temporarily grow up to the lazy threshold before a split is forced.

The amortised analysis is deliberately close to the paper's intuition.  A
lazy Insert searches the bucket directory, inserts into one bucket, and only
splits when the bucket has accumulated enough excess entries.  The potential
function records overflow beyond \(M\) together with a small slack term for
the number of buckets.  The proofs use a scaled credit invariant, so the
actual cost of a split is paid by the overflow accumulated by previous
insertions.  After the split the potential drops, and the amortised Insert
cost is the directory-search charge plus a constant amount of update work.

BatchPrepend receives a list of new candidates that belong before the
currently exposed range.  The concrete operation rebuckets that batch and
prepends the resulting buckets in one operation, which is why its accounting
is per item rather than a repeated ordinary Insert cost.  Pull removes a first
bucket when the first-bucket side conditions hold, and its scanned work is
bounded by the block size.  The theorem names reflect these three conclusions:
\thm{bp_insert_cost_bound}, \thm{bp_batch_prepend_cost_bound},
and \thm{bp_pull_cost_bound}.

The session also contains a baseline sorted-list partition model used by the
exact recurrence and abstraction proofs.  Sorted lists refine the same abstract
contracts, but they search on \(N\) entries and therefore do not supply the
paper's \(N/M\) ratio.  The bucketed structure is the implementation connected
to the active partition-interface realization facts.

\section{Executable Export and Worked Example}

The final theory gives a concrete executable surface.  Vertices and edge
weights are natural numbers, graphs are finite edge lists, and the work-list
used by the executable loop is the bucketed partition from
\thy{BMSSP_Bucketed_Partition}.  The entry point
\thm{bmssp_distances} computes a finite association list of distances for
the example graph.

The executable entry point is a finite natural-weight instantiation using the
bucketed work-list.  Its functional correctness is proved independently
against the shortest-path semantics.  The asymptotic result applies to the
costed BMSSP model and bucketed operation interface; the two layers share the
same partition primitives but are distinct artifacts.  The executable layer
uses a finite upper sentinel for natural labels; the abstract correctness
layer uses the \texttt{bound = Fin real | Infinity} datatype.

The example has source \(0\) and five vertices.  The hand-computed distances
are:
\[
  d(0)=0,\quad d(1)=3,\quad d(2)=5,\quad d(3)=8,\quad d(4)=6.
\]
The Isabelle theorem
\[
  \thm{example_bmssp_correct}
\]
states that the executable function returns exactly this association list.
Its proof is \texttt{by eval}.  Thus the bundled Poly/ML evaluator executes
the generated code equations during the build; if the executable algorithm or
the expected map is wrong, the theory does not check.

The same theory contains a \texttt{value} command for human inspection and an
\cmd{export_code} command that emits SML module \texttt{BMSSP}.  The session
ROOT file includes an export directive so the generated code can be
materialized as \texttt{generated/BMSSP.ML}.

\section{Related Work and Formalization Context}

This development is carried out in Isabelle/HOL~\cite{NipkowPaulsonWenzel2002}.
Shortest-path algorithms are a classical target for mechanized verification:
Nordhoff and Lammich~\cite{NordhoffLammich2012} verified Dijkstra's algorithm
in Isabelle/HOL within the Refinement Framework, establishing functional
correctness of an executable implementation.  The present entry differs in two
respects.  First, the algorithm is the recursive bounded multi-source
subproblem of Duan, Mao, Mao, Shu, and Yin~\cite{DuanMaoMaoShuYin2025} rather
than a single global priority queue, so the correctness argument is structured
around distance ranges and a partition interface.  Second, we treat asymptotic
running time as an explicit, separately checked artifact: a costed recurrence
is proved to satisfy the \(O(m \log^{2/3} n)\) bound under an abstract
operation-cost interface, and a concrete bucketed structure is proved to
realize the primitive costs.

Verifying \emph{complexity} alongside correctness has been pushed furthest by
Haslbeck and Lammich~\cite{HaslbeckLammich2021}, whose resource-aware
refinement (the NREST monad) carries worst-case time bounds from abstract
specifications down to LLVM code.  Our cost layer is lighter weight and
deliberately first-order: it threads exact operation counts through an
inductive cost relation rather than a monad.  This choice keeps the runtime
development close to the paper's accounting, but it also makes the cost
relation sensitive to the precise shape of its recurrence rules.  We make this
sensitivity explicit as a result in its own right.  The theory
\thy{BMSSP_Cost_Totality} contains a machine-checked \emph{obstruction}: in a
sound exact-concrete loop the single step rule forces the recursive child to
reproduce the whole cumulative bounded tree, so a genuine step is possible only
at lower bound zero
(\thm{no_genuine_exact_concrete_step_at_positive_lower_bound}), and the loop
therefore admits at most one genuine step.  That theory then records, with a
fully explicit reduction, why this step rigidity makes the run-existence
hypothesis of Section~\ref{sec:headline}'s size-parametric claim
\emph{unsatisfiable} for the directed path family \(P_k\) at every \(k \ge 2\):
the one permitted step cannot, through the truncating base case, settle more
than an initial prefix of the path.  The obstruction isolates a single
over-rigid rule and is, to our knowledge, an unusually concrete worked example
of how a cost model can be faithful to per-step accounting yet fail to be total
over the very inputs it is meant to bound.  Relaxing that rule to mirror the
more permissive operational loop is identified as the design change that would
close the gap.

For the strong BMSSP variant, \thy{BMSSP_Path_Family} now contains the checked
conditional theorem
\thm{bmssp_path_family_charged_direct_insert_runtime_bigo_size_if_cost_bounded}.
It is intentionally not advertised here as a discharged path-family headline:
the theorem waits on two external facts, eventual
\thm{charged_direct_insert_top_level_cost} existence for \(P_k\) and a charged
direct-insert runtime bound.  The runtime branch reports the latter endpoint as
\thm{charged_direct_insert_top_level_cost_refined_bound_log_params_fixed_degree},
still assuming
\thm{charged_direct_insert_closed_refined_bound_log_params_fixed_degree}.  Once
those facts are proved, the document should replace the exact-concrete vacuity
discussion above by the non-vacuous charged theorem, rather than weakening the
obstruction record.

\section{Checked Source}

The remainder of this document is the checked Isabelle source.  Definitions,
theorem statements, proof scripts, and document text are all processed by the
same session build.

\input{session}

\bibliographystyle{abbrv}
\bibliography{root}

\end{document}
